\begin{proof}
\begin{enumerate}
\item \textbf{Degenerate / edge cases.}  
If $x=0$ then $\|x+y\|=\|y\|=\|0\|+\|y\|$; similarly if $y=0$. Hence the inequality holds with equality. From now on assume $x\neq0$ and $y\neq0$.

\item \textbf{Squaring the inequality.}  
The norm is non‑negative, so $\|x+y\|\le\|x\|+\|y\|$ is equivalent to
\begin{align*}
\|x+y\|^{2}\le(\|x\|+\|y\|)^{2}. \tag{1}
\end{align*}

\item \textbf{Expansion of $\|x+y\|^{2}$.}  
Using sesquilinearity of the inner product (linearity in the first argument, conjugate‑linearity in the second) we have
\begin{align*}
\|x+y\|^{2}
&=\langle x+y,\;x+y\rangle \\[2mm]
&=\langle x,x\rangle+\langle y,y\rangle+\langle x,y\rangle+\langle y,x\rangle \\[2mm]
&=\|x\|^{2}+2\operatorname{Re}\langle x,y\rangle+\|y\|^{2}.
\tag{2}
\end{align*}

\item \textbf{Cauchy–Schwarz inequality.}  
For any $x,y\in V$,
\begin{align*}
|\langle x,y\rangle|\le\|x\|\,\|y\|,
\tag{3}
\end{align*}
with equality iff $x$ and $y$ are linearly dependent. Since $\operatorname{Re}\langle x,y\rangle\le|\langle x,y\rangle|$, we obtain
\begin{align*}
\operatorname{Re}\langle x,y\rangle\le\|x\|\,\|y\|.
\tag{4}
\end{align*}

\item \textbf{Substitution into (2).}  
Inserting the bound (4) into (2) gives
\begin{align*}
\|x+y\|^{2}
&\le\|x\|^{2}+2\|x\|\,\|y\|+\|y\|^{2} \\
&=(\|x\|+\|y\|)^{2}.
\tag{5}
\end{align*}
Inequality (5) is precisely (1).

\item \textbf{Taking square roots.}  
Both sides of (5) are non‑negative, therefore
\[
\|x+y\|\le\|x\|+\|y\|.
\]

\item \textbf{Equality condition.}  
Equality in the final inequality requires equality in every intermediate step:
\begin{itemize}
\item Equality in (4) forces $|\langle x,y\rangle|=\|x\|\,\|y\|$, which by the equality case of Cauchy–Schwarz means $y=\lambda x$ for some scalar $\lambda$.
\item Moreover we need $\operatorname{Re}\langle x,y\rangle=|\langle x,y\rangle|$. If $y=\lambda x$ then $\langle x,y\rangle=\lambda\langle x,x\rangle=\lambda\|x\|^{2}$, so $\operatorname{Re}\langle x,y\rangle=\operatorname{Re}(\lambda)\|x\|^{2}$ and $|\langle x,y\rangle|=|\lambda|\|x\|^{2}$. Equality of the real part with the modulus forces $\lambda$ to be a non‑negative real number.
\end{itemize}
Consequently,
\[
\|x+y\|=\|x\|+\|y\|
\;\Longleftrightarrow\;
\exists\,\lambda\ge0\ \text{such that } y=\lambda x
\quad(\text{equivalently } x=\mu y,\ \mu\ge0).
\]

Thus the triangle inequality holds for all vectors in an inner‑product space, and the equality case has been fully characterized. \qed
\end{enumerate}
\end{proof}